\documentclass[10pt,a4paper]{amsart}
\usepackage[margin=19mm]{geometry}
\usepackage{amsmath,amssymb,mathtools}
\usepackage[scr=rsfs,cal=euler]{mathalfa}
\usepackage{xfrac}
\usepackage[unicode,hidelinks,bookmarks=false]{hyperref}
\newcommand{\ie}{i.e.\ }
\newcommand{\cf}{cf.\ }
\newcommand{\resp}{respectively\ }
\newcommand{\Subsection}{\subsection}
\providecommand{\nobreakdash}{\nobreak\hbox{-}}
\setlength{\emergencystretch}{2em}
\multlinegap0pt
\usepackage{amssymb}
\newtheorem{prop}{Proposition}
\DeclareMathOperator{\GL}{GL}
\DeclareMathOperator{\PGL}{PGL}
\DeclareMathOperator{\charr}{char}
\title{Sha-rigidity of adjoint Chevalley groups of types $\mathbf A_1$, $\mathbf A_2$, $\mathbf B_2$, $\mathbf G_2$ over commutative rings}
\author{Elena Bunina, Vazgen Kirakosyan, Rachel Treskunov}
\date{Source: arXiv:2604.07243v3 (CC BY 4.0)}
\begin{document}
\maketitle

\noindent\emph{Source and licence.} Sha-rigidity of adjoint Chevalley groups of types $\mathbf A_1$, $\mathbf A_2$, $\mathbf B_2$, $\mathbf G_2$ over commutative rings. Version arXiv:2604.07243v3 (CC BY 4.0), distributed by arXiv under the Creative Commons Attribution 4.0 International licence, http://creativecommons.org/licenses/by/4.0/, as recorded in the arXiv licence metadata for this exact version and shown on the abstract page https://arxiv.org/abs/2604.07243v3. Full licence notice: \texttt{licenses/arxiv-2604.07243-CC-BY-4.0.txt}.\par\smallskip

\noindent\textit{Editorial scope.} Counted unit: the article's prop prop (source0.tex lines 1383--1395). The article's complete original proof of it is delivered with it (source0.tex lines 1397--1430, one \texttt{proof} environment of 34 lines). No mathematics is written, repaired or re-derived: the statement and the proof are the article's own lines, in the article's own notation, and no retained line is substituted or re-flowed.  No further source-local context is needed: every cross-reference of every retained line resolves inside the two delivered spans.  Nothing was omitted from any retained span, so the package carries no omission record.  No retained line contains a citation, so the wrapper carries no bibliography and no bibliography entry.  No retained line contains a figure environment, an image-inclusion command or drawing code, so no image is delivered and the package declares no asset dependency.  Editorial formatting only, in addition: an AMS \texttt{amsart} preamble that restates the article's own declarations and macros used by the retained lines, namely the environments \texttt{prop} on separate counters, together with the standard packages those lines need.  The retained environments print in this document as Proposition 1 in order of appearance, whereas the article prints the same items with the numbering of its own full document.  The complete native source of the article as distributed in the arXiv e-print for this exact version is bundled unchanged as \texttt{sources/198080/source0.tex}.
\begin{prop}
Let $R$ be a commutative ring and let $G_{\mathrm{ad}}(A_2,R)$ be the adjoint Chevalley group of type~$A_2$ over $R$.
Set $\gamma=\alpha_1+\alpha_2$ and consider
\[
Y \;=\; x_{\alpha_1}(1)\,w_\gamma(1)\,x_{\alpha_2}(1)\in G_{\mathrm{ad}}(A_2,R).
\]
Let $\varphi$ be the automorphism of $G_{\mathrm{ad}}(A_2,R)$ such that
\[
\varphi(x_{\pm \alpha_1}(1))=x_{\pm \alpha_2}(1),\qquad
\varphi(x_{\pm \alpha_2}(1))=x_{\pm \alpha_1}(1).
\]
Then, if $R$ has at least one localization by a maximal ideal with residue field of characteristic $\ne 7$, the elements $Y$ and $\varphi(Y)$ are not conjugate in $G_{\mathrm{ad}}(A_2,R)$.
\end{prop}

\begin{proof}
Assume that $Y$ and $\varphi(Y)$ are conjugate over $R$.
Then they remain conjugate over every localization $R_{\mathfrak m}$, and hence also over the residue field of every such localization.
Choose a maximal ideal $\mathfrak m$ for which the residue field has characteristic different from~$7$.
Reducing to that residue field and then passing to an algebraic closure, we may work inside $\PGL_3(K)$ for an algebraically closed field $K$ with $\charr K\ne 7$.
In the standard realization,
\[
w_\gamma(1)=
\begin{pmatrix}
0 & 0 & 1\\
0 & 1 & 0\\
-1& 0 & 0
\end{pmatrix},
\qquad
Y=
\begin{pmatrix}
0 & 1 & 2\\
0 & 1 & 1\\
-1& 0 & 0
\end{pmatrix}.
\]
Since the graph automorphism sends $x_\gamma(1)$ to $x_\gamma(-1)$, it sends $w_\gamma(1)$ to $w_\gamma(-1)$, and therefore $\varphi(Y)$ lifts to
\[
Y'=
\begin{pmatrix}
0 & 0 & -1\\
1 & 2 & 0\\
1 & 1 & 0
\end{pmatrix}.
\]
If $Y$ and $Y'$ were conjugate in $\PGL_3(K)$, then there would exist $g\in \GL_3(K)$ and $\lambda\in K^*$ such that $gYg^{-1}=\lambda Y'$.
Comparing determinants gives $\lambda^3=1$, while comparing traces gives $1=2\lambda$, so $\lambda=1/2$.
Hence $2^3=1$ in $K$, i.e. $7=0$, a contradiction.
\end{proof}

\end{document}
